\paragraph{About this part.} The last part looks across the whole survey. It collects the open problems that the
previous parts identified, grouped by the part they come from (Section~\ref{sec:agenda}), and states the survey's
conclusion (Section~\ref{sec:conclusion}).

% ================================================================
\section{Research agenda}
\label{sec:agenda}

The survey has identified places where the literature stops short of what a trader needs. They are grouped below
by the part of the survey they come from.

\subsection{Models}

The open problems in this group concern what the models of Part~II predict and how they are built.

\paragraph{Execution-aware objectives.} Most prediction models stop at
an estimate of $\P(\Delta m_{t,H} > 0 \mid X_t)$. The more useful objective is an estimate of
$\E[\Pi \mid a_t, X_t]$, the expected profit of each available action,
combining prediction, fill probability, queue position and adverse selection. Calibrated
confidence~\citep{manoharan2026uqlob} is a step toward it: it tells an execution policy when to act.

\paragraph{Message-level foundation models for execution.} Pretrained message-level encoders have
so far been evaluated mainly on mid-price and next-message prediction~\citep{linna2025lobert}.
Fine-tuning them for fill probability and conditional mark-outs, evaluated in queue-exact replay,
is an open direction.

\paragraph{State-space and long-context models for event streams.} State-space sequence models are well
suited to long, irregular event streams, yet apart from one hardware-oriented design~\citep{popov2026ssm} we found no established LOB forecaster built on them.
Testing them against simple baselines on the same inputs---not against each other---is a clear
next step.

\paragraph{Near-criticality and learned models.} Estimated branching ratios of order flow are high,
from about $0.7$--$0.8$ to about one depending on method. How deep sequence
models behave on near-critical, long-memory inputs, and whether Hawkes structure built into them
(Section~\ref{sec:hybrid}) improves robustness, is untested.

\paragraph{A parsimony audit.} For each widely cited deep LOB model, re-estimate its incremental
economic value and deflated performance against the best low-capacity baseline on the same data and
fill model. Such a meta-study would directly test the central claim of this survey.

\subsection{Execution and simulation}

The open problems in this group concern how predictions become orders and how simulators measure the result
(Sections~\ref{sec:sim-needs} and~\ref{sec:simulators}).

\paragraph{Queue-aware benchmarks.} Queue position should be a first-class state variable.
Benchmarks should report performance conditional on queue rank, residual volume, spread and expected
depletion time, and should measure the size of the queue-rank effect on fill probability and on
conditional mark-outs across assets and regimes, rather than take a single model estimate as
given.

\paragraph{End-to-end signal and execution.} Most models output a forecast and leave execution to a
separate rule. Models whose output is the decision itself---quote, keep, cancel, cross---and whose
training objective is execution value would close the gap that this survey describes. Reinforcement
learning in LOB simulators~\citep{lalor2025nhp,frey2023} is one route, ideally validated in queue-exact replay; its results must be
judged against model-free benchmarks in the same simulator.

\paragraph{Causal rather than predictive questions.} The questions that matter for execution are
causal: they ask what happens to the queue and the price \emph{if} an order is placed at a given price. Replay with a
no-impact assumption answers this only for small orders. Models of the market's response to one's own
orders, and ways to validate them without live experiments, remain open.

\paragraph{Generative evaluation.} Generators should be evaluated by market invariants and downstream
utility rather than by visual or marginal similarity alone~\citep{lobbench2025}.

\paragraph{Cross-asset information after costs.} An open question is when cross-asset information remains
valuable after costs, latency, multiple testing and parameter instability.

\paragraph{Simulators against live fills.} The simulators of Table~\ref{tab:simulators} have not been compared on the
same orders. Running one strategy through several of them and through paper trading, and comparing fill rates and
post-fill mark-outs with the live fills, would measure how much each simulator's assumptions---estimated queue
position, simulated order flow, no market impact---cost in accuracy.

\paragraph{Replay with a reacting market.} Replay cannot measure the market's response to one's own orders; a
model-based simulator can, mechanically (Section~\ref{sec:kaspar-fills}). Combining the two---recorded flow for the
rest of the market, a fitted model for the response to the strategy's orders---and validating the response against
live trading is open.

\subsection{Hardware}

The open problems in this group follow from Section~\ref{sec:hardware}.

\begin{itemize}[leftmargin=1.5em,itemsep=1pt]
\item Training objectives that penalise expected inference latency and resource use on a target
  device, so that models trade accuracy against deployability explicitly.
\item Which components of spatial-temporal LOB models survive mapping to programmable logic without
  losing their economic value.
\item Hawkes intensities in hardware: exponential kernels update in constant time and suit pipelines,
  and a power-law kernel can be approximated by a few exponentials, each with its own running
  sum~\citep{bochud2007}; how many are needed at the precision of fixed-point hardware, and how
  state-dependent kernels map to pipelines, are open.
\item Whether the statistical properties of a strategy developed on one hardware and latency profile
  survive deployment on another---the research--production gap of Section~\ref{sec:kaspar-determinism}, at the hardware layer.
\item Public, auditable benchmarks that cover book building and inference together: STAC-T1 audits a whole
  tick-to-trade path with simple logic, and STAC-ML audits inference alone, excluding data ingestion and features;
  neither evaluates a learned model on a book it has built.
\end{itemize}

% ================================================================
\section{Conclusion}
\label{sec:conclusion}
% ================================================================

This section states the survey's conclusion, the evidence for it, the infrastructure needed to test it, and the order
of operations it recommends to a reader who must choose a model.

The literature on limit-order-book modelling has become increasingly sophisticated at predicting the market, while the
translation from prediction to executable value has remained comparatively underdeveloped. A model is useful not
because it predicts the next price move but if the information it extracts survives every link of the chain
\[
  \text{LOB state} \to \text{representation} \to \text{prediction} \to \text{decision} \to
  \text{queue} \to \text{execution} \to \text{P\&L}
\]
to positive out-of-sample value; the quantity that matters is the expected profit of an action given the market state
and the order's place in the queue (Equation~\eqref{eq:thesis}).
\inwords{the book is turned into model inputs, the inputs into a forecast, the forecast into an order, the order
waits in a queue and may be filled, and the fill produces a profit or loss. Each link can destroy information or add
uncertainty.}
Seven conclusions follow from the evidence reviewed.
\begin{enumerate}[leftmargin=1.8em,itemsep=2pt]
\item \textbf{Representation matters.} With the network held fixed, the input decides most of the result.
\item \textbf{The prediction target matters.} How the label is built can decide the reported accuracy.
\item \textbf{Queue position matters.} Whether a passive order fills, and how good the fill is, depends on its place
  in the queue.
\item \textbf{Execution matters.} Resting and crossing can cost the same once adverse selection is counted.
\item \textbf{Simulation assumptions matter.} A fill rule that ignores the queue overstates results; a replay without
  market impact answers questions about small orders only.
\item \textbf{Latency matters.} Cost rises with delay. A slow model on the hot path misses signals, leaves passive
  orders to be picked off and lets message queues grow; a model's evaluation time is part of its specification.
\item \textbf{Complexity must earn its cost.} A more complex model is justified only by incremental out-of-sample
  economic value.
\end{enumerate}

Order-flow imbalance, queue imbalance, Hawkes intensities and first-passage models give
interpretable low-dimensional descriptions of market dynamics. Deep and pretrained models give
representational capacity and, increasingly, calibrated confidence; most of them are learned analogues of
the classical signal models rather than of the mechanistic ones (Section~\ref{sec:deep-adds}). Meta-models and
structural hybrids combine the two, either by feeding classical outputs to a learner or by building classical
structure inside a network. Much of the accuracy reported for deep models depends on how the label and the data
split are built: with a label that averages past prices, a re-test of DeepLOB on Apple was right 65.9\% of the time
and 54.6\% with a label measured from the current price~\citep{lee2024predictability}, and in our synthetic example a
network reached 57.6\% on a market with nothing to predict (Section~\ref{sec:deep-train}). The choice of input
matters as much: with the network fixed, order-flow inputs were among the best in 89\% of cases and raw snapshots in
11\%~\citep{lucchese2024}. The conclusion is not that simple models are superior. It is that complexity is an
empirical hypothesis: a more complex model is justified when it delivers incremental out-of-sample
economic value that survives realistic execution, temporal validation, and the multiple-testing
burden of its own selection.

Two kinds of infrastructure make that standard testable. The first is the simulator, and Kaspar is the survey's
instrument for it: a queue-exact replay of the exchange's own feed, deterministic, with separate latencies, running
many days at once and the same code in replay, paper and live trading (Section~\ref{sec:kaspar}). On a year of
E-mini data it shows the conclusions above in numbers: resting and crossing cost the same, about $+0.095$ ticks per
contract, post-fill mark-outs show adverse selection completed within a second, and cost rises steadily with
latency~\citep{maciejewski_shadow}. Agent-based, model-based and
generative simulators let the market react to the researcher's orders, but their queues are only as realistic as the
simulated order flow; a replay keeps the real market but cannot react. Exact queue position in the recorded market
needs order-level data and an order-level book; among the replay tools compared, hftbacktest (with order-level
data), lobsim and Kaspar provide it without altering the recorded market (Section~\ref{sec:sim-survey}). In a replay built from actors, a model of
Part~II becomes one more actor that reads the book and sends its predictions to the strategy, and the same code runs
in replay, paper and live trading (Section~\ref{sec:kaspar-actor}).

The second is hardware. Book building takes about 8\,\textmu s from socket to book in Kaspar's software path and
hundreds of nanoseconds or less in published FPGA designs; published inference times of order-book models range from
about 25\,\textmu s to about 23\,ms, about half of them above one millisecond (Table~\ref{tab:lob-latency}). Section~\ref{sec:hw-why}
identifies four causes---recurrence that must run step by step, attention whose cost grows with the square of the
window, memory traffic at a batch of one, and fixed overheads per operation---and Section~\ref{sec:hw-what} lists the
candidate remedies: shorter or incremental windows, smaller and lower-precision models, fused kernels, and moving the
regular stages to FPGAs. All of these inference times exceed the CME publisher period of about $7.5\,\mu$s, the
time within which a stage on the hot path must handle each packet to avoid queueing~\citep{maciejewski2026packets}.
A model too slow for the hot path does not just answer late: it misses signals, leaves passive orders to be picked
off, and lets the message queue in front of it grow. Hardware decides how much of a signal's value survives to the
moment of execution, but it cannot rescue a signal whose life is shorter than the pipeline that computes it.

For a reader who must choose a model, the survey's guidance reduces to an order of operations.
Decide first what the output is for---aggressive trading, passive execution, market making or
simulation---and therefore which of the targets of Section~\ref{sec:deep} matters. Establish the
strongest simple baseline on the best available representation: order-flow imbalance, queue
imbalance, the micro-price, the queue race, and for passive orders the volume ahead. Evaluate it with
proper scores and economic metrics in a queue-exact replay, counting every configuration tried. Only
then add capacity---a shallow network, a sequence model, a hybrid---and keep it only if the increment
survives the same test.

A useful future literature should therefore report not only whether a model predicts the market,
but what information it extracts, what decision it changes, how that decision is executed, and
whether the incremental value survives a properly constructed null.

\begin{table}[!htbp]
\centering\small
\caption{Evidence status of the main claims in this survey.}
\label{tab:evidence}
\begin{tabularx}{\linewidth}{Xl}
\toprule
Claim & Status \\
\midrule
Short-horizon mid changes are approximately linear in OFI, with slope inversely related to
  depth~\citep{cks2014} & established \\
Hawkes stationarity when $\rho(\Gamma) < 1$ and $\bar\lambda = (I-\Gamma)^{-1}\mu$ & established \\
Input representation can matter as much as the architecture & synthesis \\
Naive ``touch'' backtests overstate passive fills relative to queue-exact replay & synthesis;
  argued from price--time priority \\
Front-of-queue positions are worth more than back-of-queue positions, with adverse selection
  increasing in queue position~\citep{moallemi2017} & established in a model; size across regimes open \\
Hybrid inputs derived from the same stream add little to a sufficiently large network & open
  hypothesis \\
Labels that average past prices inflate reported direction accuracy~\citep{lee2024predictability} & one re-test on
  real data; reproduced on synthetic data here \\
Order-flow inputs beat raw snapshots with the network held fixed~\citep{lucchese2024} & established in one study
  (10 stocks) \\
Published inference times of LOB models range from tens of microseconds to tens of milliseconds, above book-update times & measured in five
  studies under different conditions \\
Exact own-order queue position in the recorded market requires order-level data and an order-level book &
  synthesis; argued from price--time priority \\
\bottomrule
\end{tabularx}
\end{table}
